% ─── Capítulo 0: Introducción ─────────────────────────────────────────────────
\chapter{Introducción al Compilador TPP}

\section{Descripción General}

El compilador \textbf{TPP} (\textit{Tiny Programming Processor}) es un compilador completo de un lenguaje de programación imperativo de propósito educativo, diseñado con sintaxis en español. El objetivo fundamental del proyecto es traducir código fuente escrito en el lenguaje TPP a instrucciones de ensamblador \textbf{RISC-V RV32I}, que luego pueden ser ejecutadas directamente en un procesador RISC-V de ciclo único implementado en \textbf{Logisim Evolution}.

\begin{infobox}[Objetivo del Compilador TPP]
Transformar código fuente en español al lenguaje ensamblador RV32I, permitiendo su ejecución en hardware real simulado dentro de Logisim. El resultado final (\texttt{output.s}) puede cargarse directamente en la RAM del circuito para su ejecución.
\end{infobox}

\section{Arquitectura General del Compilador}

El compilador sigue el pipeline clásico de traducción dividido en dos grandes etapas:

\begin{itemize}
    \item \textbf{Frontend}: Análisis del código fuente para verificar su validez léxica, sintáctica y semántica. Produce el Árbol Sintáctico Abstracto (AST).
    \item \textbf{Backend}: Transformación del AST en código ejecutable, pasando por optimización, representación intermedia y generación final de ensamblador RV32I.
\end{itemize}

\begin{figure}[H]
\centering
\begin{tikzpicture}[node distance=0.7cm, auto,
    stage/.style={rectangle, draw=tppblue, fill=tppblue!8,
                  text width=5.5cm, text centered, rounded corners,
                  minimum height=0.9cm, font=\small\bfseries},
    file/.style={rectangle, draw=gray, fill=gray!10,
                 text width=3.5cm, text centered, dashed,
                 minimum height=0.7cm, font=\footnotesize},
    arrow/.style={thick,->,>=stealth,color=tppblue},
    label/.style={font=\small\itshape, text=gray}]
    
    % Nodos del pipeline
    \node [file]   (src)  {archivo.to};
    \node [stage]  (lex)  [below=of src]  {1. Analizador Léxico};
    \node [stage]  (par)  [below=of lex]  {2. Analizador Sintáctico};
    \node [stage]  (sem)  [below=of par]  {3. Analizador Semántico};
    \node [stage]  (opt)  [below=of sem]  {4. Optimización};
    \node [stage]  (ir)   [below=of opt]  {5. Código Intermedio};
    \node [stage]  (gen)  [below=of ir]   {6. Generación de Código};
    \node [file]   (out)  [below=of gen]  {output.s (RV32I)};
    
    % Etiquetas laterales de archivos
    \node [label, right=1cm of lex] {lexer.l};
    \node [label, right=1cm of par] {parser.y + ast.c};
    \node [label, right=1cm of sem] {semantic.c + symtab.c};
    \node [label, right=1cm of opt] {opt.c};
    \node [label, right=1cm of ir]  {ir.c};
    \node [label, right=1cm of gen] {codegen.c};
    
    % ─── Línea divisoria Frontend / Backend ───────────────────────────────
    % Dibujada ENTRE sem y opt: usando coordenadas calculadas manualmente
    % yshift=-0.6cm baja la línea para estar debajo del nodo sem
    \draw[thick, dashed, color=gray!70,
          postaction={decorate,
            decoration={markings,
              mark=at position 0.0 with {\node[above right, font=\small\bfseries, color=tppred,  xshift=0.1cm] {FRONTEND};},
              mark=at position 1.0 with {\node[above left,  font=\small\bfseries, color=tppgreen, xshift=-0.1cm] {BACKEND};}}}]
        ([xshift=-3cm, yshift=-0.6cm]sem.south) -- ([xshift=3cm, yshift=-0.6cm]sem.south);
    
    % Flechas
    \draw [arrow] (src) -- (lex);
    \draw [arrow] (lex) -- (par);
    \draw [arrow] (par) -- (sem);
    \draw [arrow] (sem) -- (opt);
    \draw [arrow] (opt) -- (ir);
    \draw [arrow] (ir) -- (gen);
    \draw [arrow] (gen) -- (out);
\end{tikzpicture}
\caption{Pipeline completo del compilador TPP}
\label{fig:pipeline}
\end{figure}

\section{Características del Lenguaje TPP}

El lenguaje TPP soporta las siguientes características:

\begin{table}[H]
\centering
\begin{tabular}{lll}
\toprule
\textbf{Categoría} & \textbf{Construcción} & \textbf{Ejemplo} \\
\midrule
Tipos de datos & Entero, Decimal & \texttt{entero x = 5;} \\
Arreglos & Estáticos de tamaño fijo & \texttt{entero arr[10];} \\
Control de flujo & si/sino, mientras, para & \texttt{si (x > 0) \{ ... \}} \\
Funciones & Con parámetros y retorno & \texttt{fun sumar(entero a) -> entero} \\
E/S & Imprimir y leer & \texttt{imprimir("Hola");} \\
Comentarios & Línea simple & \texttt{// comentario} \\
Operadores & Aritméticos y relacionales & \texttt{+, -, *, /, ==, !=, <, >} \\
\bottomrule
\end{tabular}
\caption{Características del lenguaje TPP}
\label{tab:caracteristicas}
\end{table}

\section{Arquitectura Objetivo: RISC-V RV32I}

El código generado se dirige a la ISA \textbf{RISC-V RV32I} (32-bit Integer), ejecutado en un procesador de ciclo único dentro de Logisim. Las características relevantes son:

\begin{itemize}
    \item \textbf{Registros}: 32 registros de propósito general (x0--x31).
    \item \textbf{Stack Pointer}: Registro \texttt{sp} (x2) para manejo del stack.
    \item \textbf{Registros temporales}: \texttt{t0--t6} para evaluación de expresiones.
    \item \textbf{Registros de argumentos}: \texttt{a0--a7} para paso de parámetros.
    \item \textbf{TTY I/O}: La dirección \texttt{0xff000004} es el puerto de salida para imprimir.
    \item \textbf{Multiplicación}: Implementada por software (\texttt{\_\_soft\_mul}) ya que RV32I base no incluye \texttt{mul}.
\end{itemize}

\section{Estructura de Archivos del Proyecto}

\begin{table}[H]
\centering
\begin{tabular}{lll}
\toprule
\textbf{Archivo} & \textbf{Líneas} & \textbf{Función} \\
\midrule
\texttt{lexer.l} & 85 & Análisis léxico (flex) \\
\texttt{parser.y} & 321 & Análisis sintáctico (bison) + orquestación \\
\texttt{ast.c / ast.h} & 488 / 96 & Árbol Sintáctico Abstracto \\
\texttt{semantic.c} & 323 & Análisis semántico y tabla de símbolos \\
\texttt{symtab.c / symtab.h} & 141 / 30 & Tabla de símbolos \\
\texttt{opt.c} & 298 & Optimizaciones del AST \\
\texttt{ir.c} & 178 & Código intermedio (TAC) \\
\texttt{codegen.c} & 457 & Generación de código RV32I \\
\bottomrule
\end{tabular}
\caption{Archivos del compilador TPP}
\label{tab:archivos}
\end{table}
